\section{Richardson Extrapolation}

Refinement is not automatically improvement. A finer trial may expose a
residue that the coarser trial could not see. It may also spend more cost
without narrowing the reading that matters. The grammar therefore cannot say
``run deeper'' and call the result knowledge. It must specify which deeper
trial refines the old one, how the old residue is carried, and what error is
being cancelled or bounded.

The Richardson-style move begins with two finite depths. At depth \(d\), the
trial reads an aperture value \(a_d\). At doubled depth \(2d\), it reads
another value \(a_{2d}\). Both readings are finite. Both are bounded. The
doubled reading is not a completed limit. It is a second trial that must
preserve the seed, residue, and aperture language of the first trial while
sampling the bracket more tightly.

The grammar licenses extrapolation only when the leading sampling error can
be read with a sign. A coarse depth may read the value from one side of the
aperture. The doubled depth may force the leading error to the other side.
When that happens, the disagreement between \(a_d\) and \(a_{2d}\) is not
mere noise. It is a signed witness. The error has faced the reader, crossed
the bracket, and become eligible for cancellation.

One schematic form is enough:
\[
  a_d = \alpha + e_d,
  \qquad
  a_{2d} = \alpha - \lambda e_d + r_d,
  \qquad
  0 \leq |r_d| \leq \epsilon_d .
\]
The notation records the grammar rather than founding it. The first reading
carries a leading signed error \(e_d\). The doubled reading carries the same
error with opposite face and controlled scale \(\lambda\), plus a remainder
\(r_d\). The bound \(\epsilon_d\) says how much unread residue remains after
the paired-depth comparison.

From those two finite readings the grammar may form the improved estimate
\[
  \widehat{\alpha}_{2d}
    = \frac{a_{2d}+\lambda a_d}{1+\lambda}.
\]
If the sign relation holds, the leading error cancels and the remaining
uncertainty is bounded by \(\epsilon_d/(1+\lambda)\). The chapter does not
need the formula as a ritual. It needs the obligation the formula records:
the extrapolated value is lawful only because two finite depths have trapped
the leading error between them and left a smaller explicit residue.

If the error does not change side, the extrapolation is not forced. Two
readings on the same side of the aperture may show drift, but they do not
authorize cancellation. The grammar may carry both readings, widen the
residue, or require another depth. It may not average them and pretend the
leading error has been removed. A signless disagreement is a warning, not a
Richardson certificate.

The doubling itself is a grammatical act. It is not the fantasy that every
intermediate point has been sampled. It is a finite rule: take the current
depth, refine by the agreed factor, and preserve the prior commitments unless
a new bounded residue explicitly revises them. If the doubled trial changes
the seed discipline, the aperture face, or the meaning of value, it is not a
Richardson partner. It is a different experiment.

The coastline example shows why this caution matters. A coast measured with
a shorter ruler may report a longer path because bays and inlets that were
previously bridged now enter the path. Refinement has changed the reading.
That change is not a mistake, but it also is not convergence by itself. To
extrapolate a boundary reading, the grammar needs a residual law that says
which roughness is being exposed, how the exposure scales with depth, and
why the remaining error is bounded.

The finite-precision example shows the opposite danger. Refinement can reach
a floor where new distinctions no longer enter as new events. Additional
operations then do not sharpen the value; they merely recycle existing
residue. A bounded extrapolation must know when it has reached such a floor.
The grammar therefore carries a machine-scale residue: the smallest
distinction the trial is authorized to spend at the current depth.

For \(\alpha\), the paired-depth reading has to preserve the three-rung
aperture. The charge rung must still read loop count and orientation. The
mass rung must still read strain. The value rung must still read the
normalized ratio. A doubled depth that sharpens only one rung while leaving
the others unaccounted for cannot force an alpha reading. It may be useful
data, but it is not yet the dimensionless coupling of the three rungs.

Depth doubling also carries the seed and residue from the previous section.
The second depth does not begin as if the first had never happened. Its seed
is the successor selector of the trial, and its residue includes the
unfinished remainder carried to the doubled sampling. This is why the
extrapolated value belongs to one process. It is not the agreement of two
strangers. It is the cancellation of a signed error inside one bounded
history.

The result is a finite extrapolation, not a completed revelation. The
grammar may say: these two depths, under this aperture and carried residue,
force the reading into this smaller bracket. It may not say: all future
depths have been consumed. The bound remains part of the result. A later
trial may refine the bracket again, but it must answer to the estimate and
residue already carried.

This is the Richardson obligation in the chapter. The grammar samples,
doubles, reads a sign-flipped error, cancels the leading term, bounds the
remainder, and carries the improved value forward as a finite estimate of
\(\alpha\). The next section asks what this estimate means along a path.
Two paths may deliver the same signed displacement through the rungs. They
may still spend different costs, and the alpha reading must keep that
difference visible.
